\documentclass[11pt,a4paper]{article}

% ---------------------------------------------------------------------------
% Paquetes
% ---------------------------------------------------------------------------
\usepackage[utf8]{inputenc}
\usepackage[T1]{fontenc}
\usepackage{lmodern}
\usepackage[spanish]{babel}
\usepackage[margin=2.5cm]{geometry}
\usepackage{listings}
\usepackage{xcolor}
\usepackage{hyperref}
\usepackage{enumitem}

\hypersetup{
    colorlinks=true,
    linkcolor=blue!50!black,
    citecolor=black,
    urlcolor=blue!50!black
}

% ---------------------------------------------------------------------------
% Estilo de los listados de codigo
% ---------------------------------------------------------------------------
\definecolor{codebg}{RGB}{245,245,245}
\definecolor{codekw}{RGB}{0,0,180}
\definecolor{codecom}{RGB}{0,120,0}
\definecolor{codestr}{RGB}{170,90,0}

\lstdefinestyle{cstyle}{
    language=C,
    backgroundcolor=\color{codebg},
    basicstyle=\ttfamily\small,
    keywordstyle=\color{codekw}\bfseries,
    commentstyle=\color{codecom}\itshape,
    stringstyle=\color{codestr},
    numbers=left,
    numbersep=6pt,
    numberstyle=\tiny\color{gray},
    frame=single,
    rulecolor=\color{gray!60},
    breaklines=true,
    breakatwhitespace=false,
    showstringspaces=false,
    tabsize=2,
    captionpos=b,
    xleftmargin=1.5em,
    framexleftmargin=1.2em
}
\lstset{style=cstyle}

\title{\textbf{Resumen Técnico del Firmware}\\[4pt]
\large Sistema de Control de Acceso -- FRDM-K64F}
\author{Santiago Riverós \and Ignacio Sammartino \and Rodrigo Devesa \and Leandro Yiu}
\date{TP1 -- Sistemas Embebidos}

\begin{document}

\maketitle

\begin{abstract}
\noindent
Este documento resume, con fragmentos de código reales del firmware, las
decisiones de diseño más importantes del sistema: la separación en capas
APP/HAL/MCAL, la máquina de estados principal, el driver genérico de
display, el manejo de interrupciones y el driver del lector de banda
magnética. No es una descripción exhaustiva del código: el objetivo es
explicar el \emph{por qué} detrás de las partes menos obvias de la
implementación.
\end{abstract}

\tableofcontents

% =============================================================================
\section{Arquitectura en capas: APP $\to$ HAL $\to$ MCAL}
% =============================================================================

El firmware está organizado en tres capas con una regla estricta de
dependencia: la capa de \textbf{aplicación} (\lstinline{APP_layer}) sólo
puede llamar a la capa \textbf{HAL} (\lstinline{HAL_layer}), y ésta a su vez
sólo puede llamar al \textbf{MCAL} (\lstinline{MCAL_layer}, acceso directo a
los periféricos). La única excepción documentada es que el MCAL puede
invocar, a través de un puntero a función, un callback registrado por la
capa HAL (por ejemplo, el tick del \lstinline{SysTick} llamando a la tarea
de refresco del display).

Un caso concreto de esta regla en acción: originalmente, \lstinline{App.c}
llamaba directamente a funciones del MCAL (\lstinline{SysTick_Init},
\lstinline{NVIC_SetPriority}), lo que violaba la capa HAL. La solución fue
introducir una fachada delgada, \lstinline{scheduler.h/.c}, que no agrega
lógica propia: sólo reenvía cada llamada al MCAL correspondiente, de modo
que la capa APP nunca necesita incluir un header del MCAL.

\begin{lstlisting}[caption={Fachada HAL sobre el SysTick (\texttt{scheduler.c}): reenvío puro, sin lógica propia.}, label={lst:scheduler}]
void HAL_Scheduler_Init(void)
{
    (void)SysTick_Init();
}

bool HAL_Scheduler_AddTask(hal_task_t tarea)
{
    /* systick_callback_t y hal_task_t tienen la misma firma: la fachada
     * no le agrega semantica, solo evita que la APP conozca al MCAL. */
    return SysTick_AddCallback((systick_callback_t)tarea);
}

void HAL_Scheduler_GlobalIRQDisable(void)
{
    __disable_irq();
}

void HAL_Scheduler_GlobalIRQEnable(void)
{
    __enable_irq();
}
\end{lstlisting}

Gracias a esta fachada, \lstinline{App_Init()} puede deshabilitar las
interrupciones globales durante todo el arranque (mientras los periféricos
todavía no están configurados) y rehabilitarlas recién al final, sin que la
capa de aplicación toque directamente el core:

\begin{lstlisting}[caption={Arranque seguro en \texttt{App\_Init()}: interrupciones cortadas mientras el hardware se configura.}, label={lst:appinit}]
void App_Init(void)
{
    HAL_Scheduler_GlobalIRQDisable();

    HAL_Scheduler_Init();
    /* ... inicializacion de todos los modulos HAL ... */

    HAL_Scheduler_GlobalIRQEnable();
}
\end{lstlisting}


% =============================================================================
\section{Configuración centralizada y portabilidad}
% =============================================================================

Los valores que dependen del hardware o del ajuste fino de la aplicación
viven en archivos \lstinline{_cfg.h} dedicados (\lstinline{display_cfg.h},
\lstinline{fsm_cfg.h}), separados de la lógica que los usa. Esto permite
portar el driver de display a otra placa, u otro proyecto, tocando
\emph{un solo archivo}.

El ejemplo más claro es la macro \lstinline{DISPLAY_COMPOSE}, que define
cómo se arma la trama binaria que se envía a los shift-registers 74HC595.
El motor gráfico del driver (\lstinline{display.c}) no sabe cuántos
registros hay ni en qué byte cae cada campo: sólo calcula \emph{qué}
segmentos y selector hay que mostrar, y delega el \emph{cómo empaquetarlos}
en esta macro.

\begin{lstlisting}[caption={\texttt{DISPLAY\_COMPOSE} en \texttt{display\_cfg.h}: el único lugar que conoce el cableado real de los shift-registers.}, label={lst:compose}]
/* CABLEADO ACTUAL: dos registros; el primero en salir (U2) lleva los bits de
 * control (selector + LEDs) y el segundo (U1) los segmentos. */
#define DISPLAY_COMPOSE(frame, segs, sel, aux)                  \
    do {                                                        \
        (frame)[0] = (uint8_t)((aux) | (sel));   /* control  */ \
        (frame)[1] = (uint8_t)(segs);            /* segmentos*/ \
    } while (0)
\end{lstlisting}

De la misma forma, \lstinline{fsm_cfg.h} centraliza todos los tiempos y
umbrales de la máquina de estados (duración de mensajes, cantidad de
intentos fallidos, prioridad de interrupción del SysTick), con
\lstinline{_Static_assert} que verifican en tiempo de compilación que la
configuración sea coherente:

\begin{lstlisting}[caption={Chequeos de coherencia en tiempo de compilación (\texttt{fsm\_cfg.h}).}, label={lst:staticassert}]
#define FSM_MS_ACCESO_MENSAJE       2000U    /* "OPEn" fijo, para leerlo       */
#define FSM_MS_ACCESO_TOTAL         5000U    /* total con el pestillo abierto  */

_Static_assert(FSM_MS_ACCESO_MENSAJE <= FSM_MS_ACCESO_TOTAL,
               "OPEn fijo no puede durar mas que el total de puerta abierta");
\end{lstlisting}


% =============================================================================
\section{Máquina de estados principal}
% =============================================================================

El comportamiento de la aplicación es una FSM clásica de eventos y estados,
declarada en \lstinline{fsm.h} y despachada en \lstinline{FSM_Update()}. Los
eventos representan entradas del sistema (giros del encoder, lectura de
tarjeta, timeout de inactividad); los estados representan las pantallas /
modos en los que puede estar el sistema.

\begin{lstlisting}[caption={Enumeraciones de eventos y estados de la FSM (\texttt{fsm.h}).}, label={lst:fsmenums}]
typedef enum
{
    EV_NONE,
    EV_ENCODER_GIRO_DER, EV_ENCODER_GIRO_IZQ,
    EV_ENCODER_GIRO_PRESIONADO_DER, EV_ENCODER_GIRO_PRESIONADO_IZQ,
    EV_ENCODER_PRESIONA, EV_ENCODER_DOBLE_TOQUE,
    EV_ENCODER_MANTENER, EV_ENCODER_LIBERA_PERILLA,
    EV_CARD_OK, EV_CARD_ERROR,
    EV_TIMEOUT
} fsm_event_t;

typedef enum
{
    STATE_STANDBY, STATE_LECTURA_ID, STATE_INGRESO_PIN,
    STATE_ACCESO_OK, STATE_PIN_INCORRECTA_DELAY, STATE_BLOQUEO_ID,
    STATE_MODO_ADMIN, STATE_ERROR
} fsm_state_t;
\end{lstlisting}

El despacho es un \lstinline{switch} sobre el estado actual: cada estado es
una función \lstinline{SubFSM_*} independiente que recibe el evento y
devuelve el próximo estado. Antes de llegar a ese switch, \emph{dos}
eventos se interceptan de forma global porque aplican sin importar en qué
estado esté el sistema: el reset por mantener presionada la perilla, y el
ajuste de brillo.

\begin{lstlisting}[caption={\texttt{FSM\_Update()}: eventos globales interceptados antes del despacho por estado.}, label={lst:fsmupdate}]
void FSM_Update(fsm_event_t evento)
{
    /* RESET GLOBAL: disponible desde cualquier estado, salvo penalizado */
    if (evento == EV_ENCODER_MANTENER)
    {
        if (f_penalizado == false)
        {
            FSM_Init();
            return;
        }
    }

    /* AJUSTE DE BRILLO GLOBAL: idem, no pasa por ningun sub-estado */
    if (evento == EV_ENCODER_GIRO_PRESIONADO_DER)
    {
        HAL_Brightness_Display(INCREASE);
        return;
    }
    /* ... */

    switch (estado_actual)
    {
        case STATE_STANDBY:     estado_actual = SubFSM_StandBy(evento);     break;
        case STATE_LECTURA_ID:  estado_actual = SubFSM_Lectura_ID(evento);  break;
        /* ... resto de los estados ... */
    }
}
\end{lstlisting}


% =============================================================================
\section{Agrupación de estado relacionado en \textit{structs}}
% =============================================================================

Cada submáquina de la FSM que necesita más de una variable de contexto las
agrupa en un \lstinline{struct} propio, en vez de dejarlas como variables
sueltas del archivo. Esto deja explícito qué datos pertenecen a qué
submáquina, y evita que se puedan actualizar de forma inconsistente entre
sí (por ejemplo, tocar el cursor del menú admin sin resetear también el
flag de refresco).

\begin{lstlisting}[caption={Contexto del modo administrador agrupado en un \textit{struct} (\texttt{fsm.c}).}, label={lst:adminctx}]
typedef struct
{
    admin_substate_t substate;
    uint8_t          cursor;              /* opcion resaltada del menu     */
    uint8_t          usuario_sel;         /* indice del usuario en gestion */
    bool             f_creando_usuario;   /* alta nueva vs edicion         */
    bool             f_refrescar;         /* se acaba de (re)entrar        */
} admin_ctx_t;

static admin_ctx_t s_admin =
{
    .substate = ADMIN_VIEW_USERS, .cursor = 0U, .usuario_sel = 0U,
    .f_creando_usuario = false, .f_refrescar = true
};
\end{lstlisting}

De forma análoga, todos los estados con pantalla temporizada (acceso
concedido, error, bloqueo, ID duplicado) comparten \emph{una sola}
instancia de un cronómetro genérico, porque nunca hay dos activos al mismo
tiempo:

\begin{lstlisting}[caption={Cronómetro genérico reutilizado por todos los estados temporizados.}, label={lst:cronometro}]
typedef struct
{
    bool     f_activo;
    uint32_t inicio_ms;
    uint8_t  fase;      /* 0 = mensaje fijo, 1 = animacion (segun el estado) */
} cronometro_t;
\end{lstlisting}


% =============================================================================
\section{Temporización no bloqueante}
% =============================================================================

Todo el firmware evita \lstinline{delay()} bloqueantes: los estados
temporizados comparan la marca de tiempo guardada contra
\lstinline{HAL_GetTick()} en cada vuelta del súper-loop. Esto permite que
el sistema siga respondiendo a la interfaz (encoder, lector) mientras, por
ejemplo, se muestra el mensaje de acceso concedido.

\begin{lstlisting}[caption={Patrón de temporización no bloqueante en \texttt{SubFSM\_Acceso\_OK()}.}, label={lst:accesook}]
static fsm_state_t SubFSM_Acceso_OK(fsm_event_t evento)
{
    uint32_t transcurrido;

    if (s_cron.f_activo == false)
    {
        s_cron.f_activo  = true;
        s_cron.inicio_ms = HAL_GetTick();
        HAL_Message_Disp("OPEn");
        HAL_Buzzer_On();
    }

    transcurrido = HAL_GetTick() - s_cron.inicio_ms;

    if (transcurrido >= FSM_MS_ACCESO_TOTAL)
    {
        s_cron.f_activo = false;
        HAL_Buzzer_Off();
        FSM_Init();
        return STATE_STANDBY;
    }
    return STATE_ACCESO_OK;
}
\end{lstlisting}


% =============================================================================
\section{Ventana deslizante para buffers más largos que el display}
% =============================================================================

El display físico tiene 4 dígitos, pero el ID a ingresar tiene 8 y el PIN
hasta 5. La función \lstinline{Mostrar_VentanaIngreso} resuelve esto
dibujando un recorte de 4 posiciones del buffer lógico que se corre a
medida que el cursor avanza, de forma que el dígito en edición quede
\emph{siempre} visible en la última posición física una vez que el cursor
sale de las primeras 4.

\begin{lstlisting}[caption={Ventana deslizante: proyecta un buffer lógico largo sobre 4 dígitos físicos.}, label={lst:ventana}]
static void Mostrar_VentanaIngreso(const uint8_t *buffer, uint8_t cursor, bool f_oculto)
{
    uint8_t ancho  = HAL_Display_GetDigits();
    uint8_t offset = (cursor >= ancho) ? (uint8_t)(cursor - ancho + 1U) : 0U;
    uint8_t p;

    HAL_Blink_Digit_Disp(ALL, STOP);

    for (p = 0U; p < ancho; p++)
    {
        uint8_t logico = (uint8_t)(p + offset);

        if (logico < cursor)
        {
            /* Ya confirmado: oculto ('-') para PIN, valor real para ID. */
            HAL_Send_Digit_Disp(p, f_oculto ? (disp_char_t)'-'
                                            : Glifo_FSM(buffer[logico]));
        }
        else if (logico == cursor)
        {
            /* Cursor: siempre descubierto y titilando. */
            HAL_Send_Digit_Disp(p, Glifo_FSM(buffer[logico]));
            HAL_Blink_Digit_Disp(p, START);
        }
        else
        {
            HAL_Clear_Digit_Disp(p);
        }
    }
}
\end{lstlisting}


% =============================================================================
\section{Secciones críticas en el driver de display}
% =============================================================================

El driver de display expone su API pública desde el contexto de la
aplicación, pero el contenido de la pantalla también lo modifica
\lstinline{HAL_Refresh_Task()}, que corre dentro de la ISR del
\lstinline{SysTick}. Para que ninguna de las dos partes vea una VRAM a
medio escribir, toda función pública protege su acceso con una sección
crítica mínima basada en \lstinline{PRIMASK}:

\begin{lstlisting}[caption={Sección crítica mínima usando \texttt{PRIMASK} (\texttt{display.c}).}, label={lst:critsection}]
static inline uint32_t disp_lock(void)
{
    uint32_t primask = __get_PRIMASK();
    __disable_irq();
    return primask;
}

static inline void disp_unlock(uint32_t primask)
{
    __set_PRIMASK(primask);
}

void HAL_Clear_Digit_Disp(uint8_t target)
{
    uint32_t key = disp_lock();
    /* ... modifica vram[] y s_blink.mascara ... */
    disp_unlock(key);
}
\end{lstlisting}

Guardar y restaurar el \lstinline{PRIMASK} (en vez de simplemente
deshabilitar y rehabilitar las interrupciones) hace que estas secciones
críticas sean anidables sin efectos secundarios: si una función protegida
llama a otra función protegida, la interna no reabre las interrupciones
antes de tiempo.

Por requisito de la cátedra, el pin \lstinline{PIN_TEST_ISR} se pone en
alto mientras corre la ISR de refresco, lo que permite medir con
osciloscopio su tiempo de ejecución real:

\begin{lstlisting}[caption={Instrumentación para medir la duración de la ISR con osciloscopio.}, label={lst:testpin}]
void HAL_Refresh_Task(void)
{
#if (DISPLAY_USE_TEST_PIN != 0)
    gpioWrite(PIN_TEST_ISR, HIGH);
#endif
    /* ... bit-banging de la trama hacia los 74HC595 ... */
#if (DISPLAY_USE_TEST_PIN != 0)
    gpioWrite(PIN_TEST_ISR, LOW);
#endif
}
\end{lstlisting}


% =============================================================================
\section{Despacho genérico de interrupciones de GPIO}
% =============================================================================

El MCAL implementa un mecanismo de callbacks por pin para las interrupciones
de puerto, en vez de que cada handler conozca de antemano qué pines usa
cada driver. Una tabla de $5 \times 32$ punteros a función (un puerto, 32
pines) permite que cualquier módulo (encoder, lector de tarjetas) se
registre con \lstinline{gpioIRQ()} sin que el MCAL sepa qué hace ese
callback.

\begin{lstlisting}[caption={Tabla de callbacks y despacho genérico de IRQ de puerto (\texttt{gpio.c}).}, label={lst:gpioirq}]
static pinIrqFun_t irq_callbacks[5][32] = {0};

static void gpioIRQDispatch(uint8_t portId)
{
    uint32_t flags = port_ptrs[portId]->ISFR;
    uint8_t  i;

    for (i = 0U; i < 32U; i++)
    {
        if ((flags & (1UL << i)) != 0UL)
        {
            port_ptrs[portId]->ISFR = (1UL << i);   /* se limpia escribiendo 1 */

            if (irq_callbacks[portId][i] != 0)
            {
                irq_callbacks[portId][i]();
            }
        }
    }
}

void PORTB_IRQHandler(void) { gpioIRQDispatch(PB); }
\end{lstlisting}

Un detalle importante de la implementación es limpiar el flag de
interrupción \emph{antes} de invocar al callback: si se limpiara después,
un flanco nuevo que llegue mientras corre el callback se perdería junto
con el flag viejo.


% =============================================================================
\section{Driver del lector de banda magnética}
% =============================================================================

El driver del lector separa deliberadamente la parte \emph{urgente} de la
\emph{costosa}: la ISR de \lstinline{CLOCK} sólo guarda un bit por
interrupción (mínima y determinística), mientras que la decodificación de
la trama Track~2 completa (búsqueda de sentinelas, chequeo de paridad) se
hace en el lazo principal, nunca dentro de una ISR.

\begin{lstlisting}[caption={ISR mínima: un bit por interrupción (\texttt{lector.c}).}, label={lst:irqclk}]
void IRQ_CardClk()
{
    bitarray[currbit++] = !gpioRead(PIN_CARD_DATA); /* DATA es activo bajo */
    if (currbit >= MAX_BIT_LENGTH)
    {
        status = LECTOR_TIMEOUT;
        gpioIRQ(PIN_CARD_CLK, GPIO_IRQ_MODE_DISABLE, IRQ_CardClk);
    }
}
\end{lstlisting}

Un problema real detectado en la placa: el pin \lstinline{ENABLE} del
lector puede rebotar mecánicamente al apoyar la tarjeta, generando un pulso
cortísimo que la ISR interpreta como un swipe completo con 0--3 bits
capturados. Sin validación, esos bits (basura de una lectura anterior) se
interpretaban como una tarjeta real. La solución fue exigir una cantidad
mínima de bits capturados antes de aceptar la lectura como válida:

\begin{lstlisting}[caption={Rechazo de rebotes de \texttt{ENABLE} por cantidad insuficiente de bits (\texttt{lector.c}).}, label={lst:bounce}]
#define MIN_VALID_BITS 30
/* Una trama valida de track 2 con el ID de 8 digitos ronda los 50+ bits;
 * un rebote deja 0-3 bits. Margen generoso. */

void IRQ_CardEnable()
{
    /* ... flanco de bajada: arranca el swipe ... */
    else
    {
        /* Flanco de subida: termino el swipe. */
        gpioIRQ(PIN_CARD_CLK, GPIO_IRQ_MODE_DISABLE, IRQ_CardClk);

        if (currbit < MIN_VALID_BITS)
        {
            /* Rebote de ENABLE o ruido: se descarta aca mismo, sin pasar
             * por RECIEVED, para que nunca se vea como un intento de
             * lectura fallido en la aplicacion. */
            status = LECTOR_WAITING;
            return;
        }
        status = LECTOR_RECIEVED;
    }
}
\end{lstlisting}


% =============================================================================
\section{Prioridades de interrupción ajustadas por timing}
% =============================================================================

La ISR del \lstinline{SysTick} (que hace el \textit{bit-banging} de la
trama del display) es la más larga del sistema. Si tuviera la misma
prioridad que la IRQ del puerto del lector de tarjetas, podría retrasar la
captura de un flanco de \lstinline{CLOCK} durante el swipe y perder un
bit. La solución fue bajarle la prioridad al SysTick para que la IRQ del
lector siempre pueda interrumpirlo:

\begin{lstlisting}[caption={Ajuste de prioridad de interrupción documentado en \texttt{App\_Init()}.}, label={lst:priority}]
/* La ISR del SysTick (bit-banging del display + encoder) es la mas larga
 * del sistema y, con la misma prioridad que PORTB (0 por defecto ambas),
 * retrasaria la captura de cada bit del lector. Se le baja la prioridad
 * al SysTick para que los flancos de CLOCK la puedan interrumpir siempre. */
HAL_Scheduler_SetPriority(APP_NVIC_PRIO_SYSTICK);   /* APP_NVIC_PRIO_SYSTICK = 1 */
inicializarLector();                                /* PORTB queda en prioridad 0 (default) */
\end{lstlisting}

Esta decisión se verificó de forma empírica con osciloscopio, comparando
el pin de instrumentación de la ISR (\lstinline{PIN_TEST_ISR}) contra el
pin \lstinline{ENABLE} del lector durante un swipe real, confirmando que
ningún pulso de la ISR del SysTick bloquea una ventana donde debería
llegar un bit del lector.

\end{document}
